% Complete proofs through Durrett 3.3.21; later Section 3.3 chain remains owed.
\section{Characteristic functions and inversion}
\label{sec:annex-p15-characteristic-functions-and-inversion}

The separate source arguments are reconstructed from the held
author version \cite[pp.~125--136]{DurrettPTEAuthor2019}.
All oscillatory inversions below are limits of specified finite
integrals.

\begin{annexitem}{annex:p15.characteristic-complex-calculus}{thm:p15-characteristic-complex-calculus}{Componentwise integration, modulus bounds and exponential calculus}
\begin{proof}
The inequalities
\(|\operatorname{Re}z|,|\operatorname{Im}z|\leq|z|
\leq|\operatorname{Re}z|+|\operatorname{Im}z|\)
prove the integrability equivalence. Real linearity applied
to the two components proves complex linearity and the
conjugation identity.
Let \(w=\int f\,d\nu\). If \(w=0\), the modulus bound
is immediate. Otherwise set \(c=\overline w/|w|\).
Then \(|c|=1\) and
\[
 |w|=\operatorname{Re}(cw)
      =\int\operatorname{Re}(cf)\,d\nu
      \leq\int|f|\,d\nu.
\]
This proves exactly the norm-convexity case of Jensen's
inequality invoked in the source, without assuming an
unproved general complex Jensen theorem.
An integrable bound for the moduli bounds both components,
so the real dominated-convergence theorem proves its
complex version. Likewise absolute integrability on a
sigma-finite product permits real Fubini for both components;
recombining proves the complex interchange.

The exponential product rule follows from the real
exponential law and the trigonometric addition identities.
Differentiating its two real components along real \(x\)
gives \(d(e^{zx})/dx=ze^{zx}\).
Starting with the fundamental theorem of calculus and then
integrating by parts repeatedly gives
\[
 e^z-\sum_{k=0}^m\frac{z^k}{k!}
       =\frac{z^{m+1}}{m!}\int_0^1(1-s)^m e^{sz}\,ds.
\]
The case \(m=0\) is the fundamental theorem; integration
by parts with the polynomial factor gives each next case.
The integral has modulus at most \(e^{|z|}/(m+1)\).
Thus the remainder after degree \(m\) has modulus at most
\(e^{|z|}|z|^{m+1}/(m+1)!\).
This tends to zero by the consecutive-ratio comparison
with \(|z|/(m+2)\). The resulting exponential series is
absolutely convergent because
\(\sum |z|^k/k!=e^{|z|}\).
Finally \(|e^{iv}|=1\), and the fundamental theorem of
calculus gives
\[
 e^{iv}-1=i\int_0^v e^{is}\,ds,
\]
whose modulus is at most \(|v|\). The triangle inequality
also bounds it by two, proving the two needed estimates.
\end{proof}
\end{annexitem}

\begin{annexitem}{annex:p15.characteristic-differentiation}{thm:p15-characteristic-differentiation}{The source integral-average proof on its full measure-space domain}
\begin{proof}
First justify the measurability needed for the interchange.
On a compact subinterval of \(I\), approximate the parameter
by endpoints of finer finite partitions. The resulting
step functions \(f(t_j,s)\) and \(g(t_j,s)\) are jointly
measurable, because their sections in \(s\) are measurable.
For each \(s\), continuity in the parameter makes these
approximations converge pointwise, including the final
endpoint by assigning its exact value. Thus \(f,g\) are
jointly measurable there. A countable exhaustion by compact
subintervals gives joint measurability on \(I\times S\).
For the same reason the Riemann sums for
\(\int_J|g(t,s)|dt\), on each compact \(J\subset I\),
are measurable in \(s\) and converge for each \(s\).
Consequently
\[
 q(s)=\int_I|g(t,s)|\,dt
\]
is measurable, by an increasing compact exhaustion.

Set \(E=\{s:q(s)>0\}\).
Each set \(\{q\geq1/m\}\) has measure at most
\(m\int q\,d\nu<\infty\), and their union is \(E\).
Thus \(\nu\) restricted to \(E\) is sigma-finite,
including any points at which \(q=+\infty\).
If \(s\notin E\), the nonnegative continuous function
\(|g(\cdot,s)|\) has integral zero on \(I\), and hence
is zero at every point of \(I\). Therefore \(g(t,s)=0\)
there for all \(t\).

For \(|h|<\delta\), the fundamental theorem of calculus
on the parameter segment gives
\[
 f(y+h,s)-f(y,s)=\int_0^h g(y+\theta,s)\,d\theta.
\]
The difference vanishes outside \(E\).
On \(E\), the product with the finite Lebesgue parameter
interval is sigma-finite, and the assumed integral of
\(|g|\) is finite. Complex Fubini is therefore applicable
there. Linearity of the integrable sections and the fact
that \(g\) vanishes outside \(E\) now give
\[
 \begin{aligned}
 u(y+h)-u(y)
 &=\int_E\int_0^h g(y+\theta,s)\,d\theta\,\nu(ds)\\
 &=\int_0^h\int_E g(y+\theta,s)\,\nu(ds)\,d\theta\\
 &=\int_0^h v(y+\theta)\,d\theta.
 \end{aligned}
\]
The same equations hold for negative \(h\) with oriented
integrals. Divide by \(h\ne0\). The difference between
the resulting average and \(v(y)\) is at most
\(\sup_{|\theta|\leq|h|}|v(y+\theta)-v(y)|\), which
tends to zero by continuity. This proves the source
integral-average argument, including its arbitrary
measure-space qualification.
\end{proof}
\end{annexitem}

\begin{annexitem}{annex:p15.characteristic-algebra}{thm:p15-characteristic-algebra}{Every source algebraic property and the mixture extension}
\begin{proof}
The map \((t,x)\mapsto e^{itx}\) is continuous and has
modulus one, so its integral against a probability law
exists for every \(t\). Evaluation at zero gives one.
Conjugating under the integral gives
\(\varphi(-t)=\overline{\varphi(t)}\); the complex
modulus bound gives \(|\varphi(t)|\leq1\).
The exponential product identity yields
\[
 |\varphi(t+h)-\varphi(t)|
 \leq\int|e^{itx}(e^{ihx}-1)|\,\mu(dx)
 =\mathbb E|e^{ihX}-1|.
\]
For finite \(X\), this last integrand tends pointwise
to zero as \(h\to0\) and is bounded by two.
Bounded convergence makes the right side tend to zero
independently of \(t\), proving uniform continuity.
The affine formula follows by factoring \(e^{itb}\)
from \(e^{it(aX+b)}\); at \(a=0\) it gives the transform
of the constant \(b\).

If \(X,Y\) are independent, their joint law is the
constructed product of their real laws. The bounded
complex integrand \(e^{itx}e^{ity}\) is absolutely
integrable against this probability product.
Fubini gives
\[
 \mathbb E e^{it(X+Y)}
   =\left(\int e^{itx}\,\mu_X(dx)\right)
      \left(\int e^{ity}\,\mu_Y(dy)\right).
\]
This also proves the factorization by expanding into
the four products of real and imaginary components.

For finitely many weights, the measure
\(\mu(A)=\sum_j\lambda_j\mu_j(A)\) is a probability law,
and integration of bounded functions is linear in these
finite measures, first for indicators and simple functions
and then by bounded convergence.
For countably many weights, countable additivity of
\(\mu\) follows by interchanging the two nonnegative
series for disjoint sets. Its mass is \(\sum_j\lambda_j=1\).
For a bounded complex function, its mixture integrals
can be interchanged with the weighted sum: the sum of
the integrals of its weighted moduli is at most its
bound times \(\sum_j\lambda_j\).
Apply this to \(e^{itx}\) to obtain the countable formula.
The equal mixture of \(X,-X\) therefore has transform
\((\varphi(t)+\varphi(-t))/2=\operatorname{Re}\varphi(t)\).
Two independent copies exist on the product of the law
with itself; their difference has transform
\(\varphi(t)\varphi(-t)=|\varphi(t)|^2\).
\end{proof}
\end{annexitem}

\begin{annexitem}{annex:p15.characteristic-elementary-examples}{thm:p15-characteristic-elementary-examples}{The source sums, integrals, convolution and mixture computations}
\begin{proof}
For the fair sign, direct summation gives
\((e^{it}+e^{-it})/2=\cos t\).
For a Poisson law with \(\lambda>0\), its normalized
masses were already proved in
Theorem~\ref{thm:p15-stirling-poisson-premises}.
The absolutely convergent complex exponential series gives
\[
 \mathbb E e^{itX}
 =e^{-\lambda}\sum_{k=0}^\infty
                 \frac{(\lambda e^{it})^k}{k!}
 =\exp\{\lambda(e^{it}-1)\}.
\]
At \(\lambda=0\), the point mass convention gives one,
also the value of this formula.

For finite \(a<b\), the uniform density integrates to
one, and for \(t\ne0\) complex real-variable integration gives
\[
 \frac1{b-a}\int_a^b e^{itx}dx
                   =\frac{e^{itb}-e^{ita}}{it(b-a)}.
\]
At \(t=0\) the integral itself is one.
Substituting \(a=-c,b=c\) gives \(\sin(ct)/(ct)\).
For two independent uniforms on \((-1/2,1/2)\), the
convolution density at \(x\) is the length of the
intersection
\((-1/2,1/2)\cap(x-1/2,x+1/2)\).
Its length is \(1-|x|\) when \(|x|<1\) and zero
otherwise. Their product transform is
\[
 \left(\frac{2\sin(t/2)}t\right)^2
                       =\frac{2(1-\cos t)}{t^2}.
\]
The value one at zero follows either from the original
expectation or from the limit \(\sin v/v\to1\).
This is the source convolution proof of the triangular
case, with the density identified exactly.

For the rate-one exponential, the density integrates
to one. On \([0,R]\) its transform integral is
\((e^{(it-1)R}-1)/(it-1)\). The numerator's exponential
term tends to zero in modulus as \(R\to\infty\).
Absolute integrability is bounded by \(\int_0^\infty e^{-x}dx\);
the limit is \(1/(1-it)\).
The mixture of this law with its reflection, each
of weight \(1/2\), has density \(e^{-|x|}/2\) and transform
\[
 \frac1{2(1-it)}+\frac1{2(1+it)}=\frac1{1+t^2}.
\]
All denominators other than the specified removable
uniform and triangular ones are nonzero for real \(t\).
\end{proof}
\end{annexitem}

\begin{annexitem}{annex:p15.characteristic-normal-transform}{thm:p15-characteristic-normal-transform}{Both source Gaussian routes with their integral justifications}
\begin{proof}
First retain the real differentiation route. Evenness of
\(\phi\) makes the integral of the bounded odd function
\(\sin(tx)\phi(x)\) zero, so
\(\varphi(t)=\int\cos(tx)\phi(x)\,dx\).
For a bounded parameter neighborhood the derivative
integrand is \(-x\sin(tx)\phi(x)\), dominated by
\(|x|\phi(x)\), whose integral is finite by the already
proved Gaussian moment calculation.
Its integral is continuous in \(t\) by dominated
convergence, and its integrated modulus over the parameter
neighborhood is finite. All hypotheses of the preceding
differentiation theorem hold. Thus
\[
 \varphi'(t)=-\int x\sin(tx)\phi(x)\,dx.
\]
Since \(\phi'(x)=-x\phi(x)\), integration by parts on
\([-R,R]\), followed by the Gaussian decay of the
endpoint terms, gives
\(\varphi'(t)=-t\int\cos(tx)\phi(x)dx=-t\varphi(t)\).
Therefore the derivative of \(e^{t^2/2}\varphi(t)\)
is zero. Its value at zero is one, proving
\(\varphi(t)=e^{-t^2/2}\).

The source's completed-square route is separately valid
after establishing the complex shift. Define
\[
 f(s,x)=(2\pi)^{-1/2}e^{-(x-is)^2/2},\qquad
 J(s)=\int_{\mathbb R}f(s,x)\,dx.
\]
Fix \(0<M<\infty\). For \(|s|\leq M\),
\[
 |f(s,x)|\leq e^{M^2/2}\phi(x),\qquad
 |\partial_s f(s,x)|
       \leq(|x|+M)e^{M^2/2}\phi(x).
\]
Both bounds are integrable. The derivative is continuous
in \(s\), its integral is continuous by dominated
convergence, and the parameter-integrated bound is finite.
The differentiation theorem therefore applies at every
real \(s\). Direct differentiation gives
\(\partial_s f=-i\,\partial_x f\).
For fixed \(s\), this \(x\)-derivative is absolutely
integrable, and its integral equals
\(\lim_{R\to\infty}[f(s,R)-f(s,-R)]=0\).
Hence \(J'(s)=0\), and \(J(s)=J(0)=1\).
Completing the square now gives the actual identity
\[
 \int e^{itx}\phi(x)\,dx=e^{-t^2/2}J(t)=e^{-t^2/2}.
\]
No complex-mean probability distribution was used.
Finally the normal law of mean \(m\) and variance \(v>0\)
is the affine image \(m+\sqrt v Z\); the characteristic
affine formula proves its stated transform. At \(v=0\)
the same formula is the transform of the constant \(m\).
\end{proof}
\end{annexitem}

\begin{annexitem}{annex:p15.characteristic-dirichlet-integral}{thm:p15-characteristic-dirichlet-integral}{The complete corrected finite-strip computation}
\begin{proof}
Fix \(A>0\). Tonelli for the nonnegative modulus gives
\[
 \int_0^A\int_0^\infty e^{-xy}|\sin x|\,dy\,dx
       =\int_0^A\frac{|\sin x|}{x}\,dx<\infty.
\]
The last integral is finite because \(|\sin x|\leq x\)
near zero and the upper endpoint is finite.
Thus Fubini may be applied to \(e^{-xy}\sin x\)
on this strip. Integrating first in \(y\) gives \(S(A)\).
For fixed \(y\geq0\), elementary real integration gives
\[
 \int_0^A e^{-yx}\sin x\,dx
       =\frac{1-e^{-Ay}(\cos A+y\sin A)}{1+y^2}.
\]
All three resulting \(y\)-integrals converge absolutely:
the first is \(\int_0^\infty(1+y^2)^{-1}dy=\pi/2\),
and the other two are bounded by \(1/A\) and \(1/A^2\).
This proves the stated identity. For its sharper error bound,
write the two correction integrals as \(J_0,J_1\).
Their complex combination satisfies
\[
 |J_0+iJ_1|
  =\left|\int_0^\infty\frac{e^{-Ay}(1+iy)}{1+y^2}\,dy\right|
  \leq\int_0^\infty\frac{e^{-Ay}}{\sqrt{1+y^2}}\,dy
  \leq1/A.
\]
The real correction \(\cos A\,J_0+\sin A\,J_1\)
is a rotated real component of \(J_0+iJ_1\), so its
absolute value is at most \(1/A\).
Also \(\pi/2-\arctan A=\int_A^\infty(1+y^2)^{-1}dy\leq1/A\).
The triangle inequality gives the source's
\(|S(A)-\arctan A|\leq2/A\) as well.
In particular \(S(A)\to\pi/2\).
For \(0\leq A\leq1\), \(|S(A)|\leq A\leq1\);
for \(A\geq1\), \(|S(A)|\leq\pi/2+2\).
This gives a uniform absolute bound on the whole half-line.

For \(\theta>0\), the affine change of variable
\(x=\theta t\), together with evenness of
\(\sin(\theta t)/t\), gives
\(R(\theta,T)=2S(T\theta)\).
For negative \(\theta\) the sign reverses, and for
\(\theta=0\) the integral is zero. These give the
unified sign formula, its uniform bound and its limit.

Absolute convergence on the full half-line fails.
For each integer \(m\geq1\), on
\([m\pi+\pi/6,m\pi+5\pi/6]\) one has
\(|\sin x|\geq1/2\) and \(x\leq(m+1)\pi\).
The integral of \(|\sin x|/x\) over this interval is
at least \(1/(3(m+1))\).
These disjoint intervals yield a divergent harmonic
lower bound. This conditional convergence is fully
compatible with absolute integrability on each finite strip.
\end{proof}
\end{annexitem}

\begin{annexitem}{annex:p15.characteristic-inversion-uniqueness}{thm:p15-characteristic-inversion-uniqueness}{The oscillatory limit, endpoint masses and uniqueness}
\begin{proof}
For fixed \(T>0\), set
\(I_T=\int_{-T}^T K_{a,b}(t)\varphi(t)\,dt\).
The kernel identity gives \(|K_{a,b}(t)|\leq b-a\)
for every \(t\), including zero.
Thus the double integral against \(dt\,\mu(dx)\) on
\([-T,T]\times\mathbb R\) has modulus integral at most
\(2T(b-a)\), and complex Fubini gives
\[
 I_T=\int_{\mathbb R}\left(
            \int_{-T}^T K_{a,b}(t)e^{itx}\,dt\right)\mu(dx).
\]
For each fixed \(x\), expand the exponentials.
The imaginary terms are odd in \(t\), and integrate
to zero on the symmetric interval. The remaining
integral is
\[
 R(x-a,T)-R(x-b,T).
\]
Only the original bounded difference kernel was used
for Fubini; the two sine integrals are evaluated at
fixed \(x\), so no moment assumption on \(X\) is hidden
in this splitting.
The Dirichlet lemma bounds their difference uniformly
in \(x,T\) and gives its pointwise limit as \(2\pi\)
on \(a<x<b\), as \(\pi\) at \(x=a,b\), and as zero
outside \([a,b]\). Bounded convergence against the
probability measure, followed by division by \(2\pi\),
proves inversion and its two atom corrections.

If two probability measures have the same transform,
choose endpoints \(a<b\) that are atoms of neither.
Their atom sets have countable union, so such endpoints
form a dense set. Inversion gives equality of their
probabilities on \((a,b)\).
For each such \(b\), let eligible \(a\to-\infty\)
and use continuity from below. The two CDFs agree at
these nonatom \(b\); right continuity extends equality
to every real \(b\). CDF uniqueness gives equality of
the Borel measures.
For \(\delta_0\), the transform is one and
\(K_{-1,1}(t)=2\sin t/t\), whose nonabsolute convergence
is exactly the Dirichlet example just proved.
\end{proof}
\end{annexitem}

\begin{annexitem}{annex:p15.characteristic-inversion-corollaries}{thm:p15-characteristic-inversion-corollaries}{Symmetry and the separate normal-sum proof}
\begin{proof}
If \(\varphi\) is real, the reflected law has transform
\(\varphi(-t)=\overline{\varphi(t)}=\varphi(t)\).
Uniqueness shows that the law equals its reflection.
Conversely equality with the reflected law makes
\(\varphi(t)=\overline{\varphi(t)}\), so it is real.
For the independent centered normal pair, multiplication
of the already calculated transforms gives
\[
 e^{-v_1t^2/2}e^{-v_2t^2/2}
                   =e^{-(v_1+v_2)t^2/2}.
\]
This is the transform of the centered normal law with
the summed variance, including its zero-variance point
mass case. Uniqueness proves the law assertion.
\end{proof}
\end{annexitem}

\begin{annexitem}{annex:p15.characteristic-integrable-inversion}{thm:p15-characteristic-integrable-inversion}{A real nonnegative normalized density}
\begin{proof}
Let \(C=(2\pi)^{-1}\int|\varphi(t)|dt<\infty\).
The proposed inverse integral exists at every \(y\)
and has modulus at most \(C\).
Moreover
\[
 |f(y+h)-f(y)|
 \leq\frac1{2\pi}\int|e^{-ith}-1|\,|\varphi(t)|\,dt
                   \longrightarrow0
\]
by dominated convergence. The bound is independent
of \(y\), so \(f\) is uniformly continuous.
Conjugating the integral and using
\(\overline{\varphi(t)}=\varphi(-t)\), followed by
the affine substitution \(s=-t\), gives
\(\overline{f(y)}=f(y)\). Hence \(f\) is real.

The inversion integral now converges absolutely because
\(|K_{a,b}\varphi|\leq(b-a)|\varphi|\).
Its absolute bound yields
\[
 \mu((a,b))+\tfrac12\mu(\{a,b\})\leq C(b-a).
\]
For any fixed \(x\), applying this to \((x-\eta,x+\eta)\)
gives \(\mu(\{x\})\leq2C\eta\).
Let \(\eta\downarrow0\); the law has no atoms.
Using the kernel's integral representation, complex
Fubini is applicable with total modulus integral
at most \((b-a)\int|\varphi|\), and gives
\[
 \mu((a,b))=\int_a^b f(y)\,dy.
\]
If \(f\) were negative at any point, its continuity
would make this integral negative on a sufficiently
small interval, contradicting nonnegativity of \(\mu\).
Thus \(f\geq0\) everywhere. Monotone convergence on
\((-N,N)\) gives \(\int_{\mathbb R}f=1\).
The probability law with density \(f\) agrees with
\(\mu\) on all bounded intervals, hence on rays by
limits and the absence of atoms, and then on all Borel
sets by CDF uniqueness. This proves every asserted
density property.
\end{proof}
\end{annexitem}

\begin{annexitem}{annex:p15.characteristic-inverse-examples}{thm:p15-characteristic-inverse-examples}{The source dual inversions and the Cauchy average}
\begin{proof}
The triangular transform
\(2(1-\cos s)/s^2\), with value one at zero, is integrable:
it is bounded on \([-1,1]\) and at most \(4/s^2\)
outside that interval. Apply the integrable inversion
theorem to its known triangular law.
Both its known density and the inverse density are
continuous, so their almost-everywhere equality is
equality at every point: a nonzero continuous
difference would stay nonzero on an interval of
positive Lebesgue measure. Therefore
\[
 \frac1{2\pi}\int_{\mathbb R}
       \frac{2(1-\cos s)}{s^2}e^{-isy}\,ds=(1-|y|)^+
                       \quad(y\in\mathbb R).
\]
At \(y=0\), this proves that the nonnegative \(f_P\)
integrates to one. Its assigned value at zero is its
continuous limit, since
\((1-\cos s)/s^2=\tfrac12(\sin(s/2)/(s/2))^2\to1/2\).
Setting \(y=-t\) now proves its stated characteristic
function.

The bilateral exponential transform \(1/(1+s^2)\)
is integrable by the arctangent antiderivative.
The same theorem, and continuity of its known density,
give at every real \(y\)
\[
 \frac1{2\pi}\int_{\mathbb R}
           \frac{e^{-isy}}{1+s^2}\,ds=\tfrac12e^{-|y|}.
\]
Set \(y=-t\) and multiply by two to obtain the
characteristic function of the normalized Cauchy
density. Finally, for an independent Cauchy sample,
the transform of its average is
\((e^{-|t|/n})^n=e^{-|t|}\).
Uniqueness identifies its law with the original law,
providing this source proof separately from direct
real convolution.
\end{proof}
\end{annexitem}

\begin{annexitem}{annex:p15.characteristic-atom-inversion}{thm:p15-characteristic-atom-inversion}{Averaged kernels, lattice orthogonality and the squared-atom sum}
\begin{proof}
For finite \(T\), absolute integrability of the bounded
exponential on \([-T,T]\) times a probability space
permits Fubini. The averaged integral is
\[
 \frac1{2T}\int_{-T}^T e^{-ita}\varphi(t)\,dt
       =\int\operatorname{sinc}(T(x-a))\,\mu(dx).
\]
The integrand has modulus at most one and tends to
\(\mathbf1_{\{a\}}(x)\). Bounded convergence proves
the atom formula.

If \(x,y\in b+h\mathbb Z\), their difference is an
integer multiple of \(h\). Direct integration shows
\[
 \frac h{2\pi}\int_{-\pi/h}^{\pi/h}e^{it(y-x)}\,dt
 =\begin{cases}1,&y=x,\\0,&y\ne x.\end{cases}
\]
Fubini against the law supported on this lattice
therefore gives the point-mass inversion.
On that support,
\(e^{i(2\pi/h)y}=e^{i(2\pi/h)b}\), proving the
stated phase-periodicity formula.

For independent copies \(X,X'\), the difference has
transform \(|\varphi|^2\).
Apply the atom formula at zero.
The diagonal in \(\mathbb R^2\) is Borel, so Tonelli
for the product probability yields
\[
 \mathbb P(X=X')=\int\mu(\{x\})\,\mu(dx)
              =\sum_{x:\mu(\{x\})>0}\mu(\{x\})^2.
\]
The integrand \(x\mapsto\mu(\{x\})=F(x)-F(x-)\)
is Borel and supported on the countable atom set,
which justifies the last equality.
If \(\varphi(t)\to0\) at positive infinity, the
conjugation identity gives the same modulus limit
at negative infinity. For each \(\eta>0\), choose
\(R\) such that \(|\varphi(t)|^2\leq\eta\) for
\(|t|>R\). Its average on \([-T,T]\) is then at
most \(R/T+\eta\) for \(T\geq R\).
It tends to zero; the sum of squared atom masses
is therefore zero, and there are no atoms.
\end{proof}
\end{annexitem}

\begin{annexitem}{annex:p15.characteristic-continuity-theorem}{thm:p15-characteristic-continuity-theorem}{The source tightness calculation and bounded-test subsequence route}
\begin{proof}
Under weak convergence, the bounded continuous real
functions \(\cos(tx)\) and \(\sin(tx)\) have converging
integrals for every fixed \(t\). Recombine their real
and imaginary parts to obtain the first implication.

For the converse, the pointwise limit is Borel measurable
because its two components are limits of continuous functions.
Pointwise limits also preserve
\(\varphi(0)=1\), \(|\varphi(t)|\leq1\) and the
conjugation symmetry.
For \(u>0\), the double integrand \(1-e^{itx}\)
has modulus at most two on the finite parameter
interval times a probability space. Fubini and
direct integration in \(t\) give
\[
 \frac1u\int_{-u}^u(1-\varphi_n(t))\,dt
       =2\int(1-\operatorname{sinc}(ux))\,\mu_n(dx).
\]
Oddness of the imaginary component identifies the
left side with its real-part version.
Since \(|\sin v|\leq|v|\), the right integrand is
nonnegative. On \(|ux|\geq2\),
\[
 2(1-\operatorname{sinc}(ux))
                    \geq2(1-|ux|^{-1})\geq1.
\]
Discarding the complementary nonnegative contribution
proves the exact displayed tail bound.

For each fixed \(u>0\), bounded convergence in the
finite \(t\)-interval gives convergence of this
left integral to
\[
 A(u)=\frac1u\int_{-u}^u
                  (1-\operatorname{Re}\varphi(t))\,dt.
\]
Continuity at zero and \(\varphi(0)=1\) imply
\(0\leq A(u)\leq2\sup_{|t|\leq u}|1-\varphi(t)|\to0\).
For a required tolerance, choose \(u\) so that \(A(u)\)
is smaller than half that tolerance. The approximating
integrals are eventually smaller than the full tolerance.
The displayed tail estimate gives the source tightness
condition; the finitely many initial probability laws
can be included by enlarging the radius, as proved in
the uniform-tail formulation of tightness.

Helly selection and tightness provide at least one
subsequence converging weakly to a proper law \(\mu\).
By the first implication its transform is \(\varphi\).
Every subsequence of the original sequence has a
further proper weak limit with that same transform;
uniqueness proves that each such limit is \(\mu\).
To retain the source's last argument explicitly, fix
a bounded continuous real \(f\). Every subsequence of
\(\int f\,d\mu_n\) has a further subsequence tending
to \(\int f\,d\mu\).
If the full scalar sequence did not tend to that
number, some positive tolerance and subsequence
would stay outside its tolerance interval, making
such a further convergence impossible.
Thus all these test integrals converge, and the
bounded-continuous characterization proves
\(\mu_n\Rightarrow\mu\).
This also proves existence and uniqueness of the
limiting law; neither is assumed from the pointwise
candidate function alone.
\end{proof}
\end{annexitem}

\begin{annexitem}{annex:p15.characteristic-moment-derivatives}{thm:p15-characteristic-moment-derivatives}{Every differentiated integral and continuity assertion}
\begin{proof}
For \(0\leq j\leq n\), \(|x|^j\leq1+|x|^n\);
thus all the relevant lower moments are finite.
Suppose the formula has been proved at order \(j-1\).
Use the differentiation theorem with
\[
 f(t,x)=(ix)^{j-1}e^{itx},\qquad
 g(t,x)=(ix)^j e^{itx}.
\]
For each fixed \(x\) the parameter function is
continuously differentiable. Its sections and their
derivatives are integrable by the moment bounds.
The derivative integral is continuous by dominated
convergence with bound \(|x|^j\), and on any parameter
interval of length \(2\delta\) its modulus integral
is at most \(2\delta\mathbb E|X|^j<\infty\).
Consequently the derivative formula holds at the
next order at every real parameter.
Starting at \(j=1\) and iterating proves all formulas;
the same dominated-convergence argument proves the
continuity of each derivative. Evaluation at zero
and complex linearity give the moment identity.
\end{proof}
\end{annexitem}

\begin{annexitem}{annex:p15.characteristic-taylor-remainder}{thm:p15-characteristic-taylor-remainder}{The source two integrations by parts, including zero and negative arguments}
\begin{proof}
Put \(I_n(x)=\int_0^x(x-s)^n e^{is}\,ds\).
For \(n=0\), \(I_0(x)=(e^{ix}-1)/i\).
For \(n\geq1\), integration by parts yields
\[
 I_{n-1}(x)=x^n/n+iI_n(x)/n.
\]
Induction in this identity proves
\[
 e^{ix}-\sum_{m=0}^n(ix)^m/m!
                         =i^{n+1}I_n(x)/n!.
\]
It is valid for negative \(x\) with oriented integrals.
Bounding the modulus of the integral by the integral
of \(|x-s|^n\) over the segment between zero and \(x\)
gives the first bound \(|x|^{n+1}/(n+1)!\).

For \(n\geq1\), subtract
\(x^n/n=\int_0^x(x-s)^{n-1}ds\) in the recurrence.
The exact same remainder can then be written
\[
 \frac{i^n}{(n-1)!}
       \int_0^x(x-s)^{n-1}(e^{is}-1)\,ds.
\]
Use \(|e^{is}-1|\leq2\) to bound its modulus by
\(2|x|^n/n!\).
For \(n=0\), this second bound is directly
\(|e^{ix}-1|\leq2\).
The first bound at \(n=0\) is \(|e^{ix}-1|\leq|x|\),
already proved. At \(x=0\) the remainder is zero
and the stated conventions give the same bound.
Taking the smaller of the two bounds proves the
full source lemma on its entire claimed domain.
\end{proof}
\end{annexitem}

\begin{annexitem}{annex:p15.characteristic-moment-expansion}{thm:p15-characteristic-moment-expansion}{The integrable minimum bound without an extra moment}
\begin{proof}
Each polynomial term is integrable under the
stated moment hypothesis. Complex linearity and
the modulus inequality followed by the Taylor
lemma at \(x=tX\) give exactly the expected
minimum bound. Its finiteness follows from the
second term, even if the expectation of the
first term alone is infinite.
Dropping the factorial denominators, each at least
one, gives the source's additional bound (3.3.3).
For \(t\ne0\), divide that bound by \(|t|^n\).
The result is
\[
 \mathbb E\min\left\{
       \frac{|t|\,|X|^{n+1}}{(n+1)!},
                         \frac{2|X|^n}{n!}\right\}.
\]
The integrand tends pointwise to zero because
\(X\) is finite, and is dominated by the
integrable second term. Dominated convergence
therefore proves \(o(|t|^n)\).
For \(n=0\) the same calculation uses the
constant dominating function two.
Taking \(n=2\) gives source Theorem~3.3.20;
its proof has used only \(\mathbb EX^2<\infty\),
and no third-moment integral is separated from
the minimum.
In the source's denominator-free form the bound is
\(t^2\mathbb E\min\{|t|\,|X|^3,2|X|^2\}\);
its integral after division by \(t^2\) tends to
zero under precisely the same domination by \(2|X|^2\).
\end{proof}
\end{annexitem}

\begin{annexitem}{annex:p15.characteristic-second-moment}{thm:p15-characteristic-second-moment}{The full Fatou converse and its zero-law consequence}
\begin{proof}
Conjugation symmetry gives
\[
 Q(h)=-\int\frac{2(1-\cos(hx))}{h^2}\,\mu(dx)\leq0.
\]
Thus it is real. If its upper limit as \(h\downarrow0\)
is greater than \(-\infty\), then the lower limit of
the nonnegative integrals on the right, without the
minus sign, is finite. Choose \(h_j\downarrow0\)
along which those integrals tend to that lower limit.
The integrands tend pointwise to \(x^2\), since
\(2(1-\cos v)/v^2\to1\).
Fatou gives
\[
 \mathbb EX^2\leq
 \liminf_j\int\frac{2(1-\cos(h_jx))}{h_j^2}\,\mu(dx)
 =-\limsup_{h\downarrow0}Q(h)<\infty.
\]
The sequence can be chosen strictly decreasing
by the definition of the lower limit on shrinking
positive intervals.
Conversely, the finite-second-moment expansion
gives \(Q(h)=-\mathbb EX^2+o(1)\).
This proves the equivalence and the full two-sided
limit.

If \((\varphi(t)-1)/t^2\to c\in\mathbb C\) as \(t\downarrow0\),
then \(Q(h)\to2\operatorname{Re}c\) along positive
\(h\), so the first part proves a finite second
moment. That also gives a finite first absolute
moment. The expansion implies
\[
 \frac{\varphi(t)-1}{t^2}
       =\frac{i\mathbb EX}{t}-\tfrac12\mathbb EX^2+o(1).
\]
A finite limit forces \(\mathbb EX=0\), and then
\(c=-\mathbb EX^2/2\) is real.
In the zero-remainder case the second moment
is zero. Nonnegativity of \(X^2\), or Markov's
inequality at every positive threshold, gives
\(X=0\) almost surely, and its transform is
identically one.
\end{proof}
\end{annexitem}

\begin{annexitem}{annex:p15.characteristic-boundaries}{thm:p15-characteristic-boundaries}{Failure at zero, nonintegrable transforms and independence}
\begin{proof}
The normal transform formula gives \(e^{-nt^2/2}\).
Its pointwise limit is one at zero and zero elsewhere.
The normal CDFs are \(\Phi(x/\sqrt n)\), tending
to \(\Phi(0)=1/2\) at every real \(x\).
Their probabilities in any fixed bounded interval
tend to zero, so tightness fails. This explains
the limiting continuity hypothesis by a direct
calculation, without any CLT.

For a uniform law on \((-1,1)\), the transform is
\(\sin t/t\); its absolute integral diverges by
the Dirichlet lemma's harmonic lower bound.
For the exponential transform its modulus is
\((1+t^2)^{-1/2}\), at least \(1/(\sqrt2\,t)\)
for \(t\geq1\), so its integral also diverges.
Both laws have the densities already specified.
Thus the integrability condition for inverse
density recovery is sufficient but is not a
necessary condition for existence of a density.

If \(X=Y\) is a fair sign, the transform of
the sum is \(\cos(2t)\); multiplication of
the marginal transforms would give \(\cos^2t\).
At \(t=\pi/2\) the two values are \(-1\) and zero.
Finally, for \(\alpha>2\),
\(e^{-|t|^\alpha}=1+o(t^2)\), by the real exponential
expansion at zero. If it were a characteristic
function, the zero-law result would make it
identically one. Its value at \(t=1\) is \(e^{-1}<1\),
a contradiction.
\end{proof}
\end{annexitem}
